\newpage
\phantomsection
\label{prf:finite-sum-rule}

\begin{remark*}[Return]
\hyperref[cor:finite-sum-rule]{Return to Corollary}
\end{remark*}

\begin{theorem*}[Finite Sum Rule]
If $f_1, f_2, \dots, f_n$ are functions on an interval $I$ to $\mathbb{R}$ that are differentiable at $c \in I$, then:
\[ \left( \sum_{i=1}^{n} f_i \right)'(c) = \sum_{i=1}^{n} f_i'(c) \]
\end{theorem*}

\begin{proof}
\textbf{Professional Standard Proof.}~\\
The proof proceeds by mathematical induction on $n$. For the base case $n=1$, the statement $(f_1)'(c) = f_1'(c)$ holds trivially. Assume the statement holds for some $k \in \mathbb{N}$, so $\left( \sum_{i=1}^{k} f_i \right)'(c) = \sum_{i=1}^{k} f_i'(c)$. For $n = k+1$, let $g = \sum_{i=1}^{k} f_i$. Then $\left( \sum_{i=1}^{k+1} f_i \right)'(c) = (g + f_{k+1})'(c)$. By the sum rule for derivatives, $(g + f_{k+1})'(c) = g'(c) + f_{k+1}'(c)$. Applying the inductive hypothesis yields $g'(c) + f_{k+1}'(c) = \sum_{i=1}^{k} f_i'(c) + f_{k+1}'(c) = \sum_{i=1}^{k+1} f_i'(c)$. By mathematical induction, the statement holds for all $n \in \mathbb{N}$.
\end{proof}

\begin{proof}
\textbf{Detailed Learning Proof.}~\\

\textbf{Step 1.} Establish the base case for induction. For $n=1$, the statement becomes $(f_1)'(c) = f_1'(c)$, which is an identity and therefore holds trivially.

\textbf{Step 2.} Formulate the inductive hypothesis. Assume the statement holds for some $k \in \mathbb{N}$, meaning that for any collection of $k$ functions differentiable at $c$, their sum satisfies $\left( \sum_{i=1}^{k} f_i \right)'(c) = \sum_{i=1}^{k} f_i'(c)$.

\textbf{Step 3.} Prove the inductive step for $n = k+1$. Define $g = \sum_{i=1}^{k} f_i$. Then the sum of $k+1$ functions can be written as $\sum_{i=1}^{k+1} f_i = g + f_{k+1}$, allowing us to reduce to the two-function case.

\textbf{Step 4.} Apply the sum rule for derivatives. Since $g$ and $f_{k+1}$ are both differentiable at $c$ (the former by the inductive hypothesis and the latter by assumption), the sum rule gives $(g + f_{k+1})'(c) = g'(c) + f_{k+1}'(c)$.

\textbf{Step 5.} Substitute the inductive hypothesis. By the inductive hypothesis, $g'(c) = \left( \sum_{i=1}^{k} f_i \right)'(c) = \sum_{i=1}^{k} f_i'(c)$. Therefore, $g'(c) + f_{k+1}'(c) = \sum_{i=1}^{k} f_i'(c) + f_{k+1}'(c) = \sum_{i=1}^{k+1} f_i'(c)$.

\textbf{Step 6.} Conclude by mathematical induction. Having established both the base case and the inductive step, the principle of mathematical induction guarantees that the statement holds for all $n \in \mathbb{N}$.
\end{proof}

\begin{remark*}[Proof structure]
This proof employs mathematical induction to extend the sum rule from two functions to an arbitrary finite number of functions. The strategy reduces the general case to repeated application of the two-function sum rule, leveraging the inductive structure to build from simple cases to the general result.
\end{remark*}

\begin{remark*}[Dependencies]~\\
\begin{itemize}
  \item \hyperref[def:derivative]{Definition of Derivative}
  \item \hyperref[thm:sum-rule]{Sum Rule for Derivatives}
  \item \hyperref[thm:mathematical-induction]{Principle of Mathematical Induction}
\end{itemize}
\end{remark*}